% Chapter 30: Orthographic Evolution
\chapter{Orthographic Evolution}

\chapterepigraph{%
  Scripts do not drift randomly.\\
  They drift toward efficiency, legibility, and the limits\\
  of the human hand at the speed it is asked to move.
}{Flyxion, \textit{Semantic Infrastructure}}

\sidenote{Connects to\\gesture fossil-\\ization (Ch.\,28)\\and constraint\\descent (App.\,A)}

Writing systems evolve. Over centuries and millennia, scripts transform
in ways that are not random but are governed by systematic pressures:
toward faster production, toward greater consistency, toward reduced
ambiguity. This chapter examines orthographic evolution as a constraint
descent process: scripts move toward configurations of lower constraint
violation under the pressures of use.

\section{Arabic: Cursive Optimization}

Arabic script represents one of the most sophisticated solutions to
the speed-distinctiveness tradeoff in the history of writing. Its
cursive form — in which letters connect and change shape with position
— is often presented as complex. The constraint-descent view reveals
it as elegant.

The four positional variants of each letter (isolated, initial, medial,
final) are not arbitrary complications. They are the result of optimizing
two competing constraints:
\begin{enumerate}
  \item \textbf{Connectivity:} cursive writing is faster than print
        because the pen does not leave the surface between letters.
  \item \textbf{Distinctiveness:} letters must remain distinguishable
        despite deformation from the baseline form.
\end{enumerate}

The positional variants represent the four stable attractors in
letter-shape space that satisfy both constraints simultaneously:
connected to the preceding letter, connected to the following letter,
connected to both, connected to neither.

\section{Latin: From Majuscule to Minuscule}

The Latin alphabet underwent a major transformation between the 4th and
9th centuries CE: from all-capitals (majuscule) to mixed case (minuscule).
This transition was driven by a single constraint: writing speed.

Minuscule letters are faster to produce than majuscule for three reasons:
\begin{itemize}
  \item \textbf{Fewer strokes:} many minuscules require fewer pen lifts
        than their majuscule equivalents.
  \item \textbf{Shorter travel:} ascenders and descenders guide the pen
        in a consistent vertical rhythm, reducing lateral travel.
  \item \textbf{Ligatures:} minuscule letters connect naturally into
        common clusters (st, ct, æ, œ), enabling cursive-like efficiency.
\end{itemize}

The majuscule-to-minuscule transition is a constraint descent: the
constraint violation functional $\kappa_\text{speed}$ (time per word)
decreased substantially with the new script.

\section{Greek, Coptic, and Constructed Scripts}

The Greek alphabet demonstrates the opposite of fossilization: active
construction. When the Greeks adapted the Phoenician alphabet to write
vowels (around 800 BCE), they were not merely adopting a system —
they were engineering one. The letters they repurposed (aleph → alpha,
he → epsilon, yod → iota) were chosen not for sound but for phonemic
coverage: which Phoenician letters sounded most like Greek vowels?

This is a constraint satisfaction problem: assign Phoenician letters
to Greek phonemes such that the assignment minimizes perceptual
distance between letter-name and phone, while maximizing coverage of
the Greek phonemic inventory.

The Coptic alphabet (derived from Greek to write Egyptian) similarly
involved active engineering: Greek letters were supplemented with
Demotic Egyptian characters for phonemes Greek did not possess.
The Coptic designers solved a gap-filling constraint problem: extend
the Greek basis to cover the full Coptic phonemic space with
minimum additional symbols.

\section{Constructed Scripts as Constraint Satisfaction}

Deliberately constructed scripts — the Korean hangul (1443 CE),
the Cherokee syllabary (1820s), the Canadian Aboriginal syllabics
(1840s), and many others — provide natural experiments in orthographic
design. Each was designed by one or a small number of inventors,
making the design choices explicit.

\begin{definition}{Constructed Script Optimality}{constructed-opt}
  A \emph{constructed script} is \emph{optimal} for a language $L$ if
  it minimizes total constraint violation:
  \[
    \kappa_\text{total} = w_1 \kappa_\text{completeness}
    + w_2 \kappa_\text{distinctiveness}
    + w_3 \kappa_\text{economy}
    + w_4 \kappa_\text{learnability}
  \]
  where $w_1, w_2, w_3, w_4$ are weights reflecting the community's
  priorities.
\end{definition}

Hangul is remarkable for its systematic design: consonants are based
on the shape of the mouth and tongue during articulation (a form of
phonetic iconicity), and syllable blocks are composed from consonants
and vowels according to a regular spatial grammar. This achieves
high learnability (the blocks encode phonetic structure) and high
distinctiveness (the spatial arrangement provides additional cues).

\section{Why Alphabets Drift Over Centuries}

Even well-designed scripts drift. The drift follows predictable patterns:

\begin{enumerate}
  \item \textbf{Speed simplification:} strokes that can be omitted
        without loss of distinctiveness are omitted. Features that
        require pen lift are merged into continuous curves.
  \item \textbf{Frequency adjustment:} common letters become simpler;
        rare letters remain complex because the frequency-weighted
        optimization favors common-letter efficiency.
  \item \textbf{Distinctiveness pressure:} when confusion between letters
        increases (due to simplification of one), the other is
        elaborated or the simplified form acquires a diacritic.
\end{enumerate}

This is precisely constraint descent: the script evolves by locally
reducing $\kappa_\text{total}$, with each generation of writers
making small modifications that persist if they reduce the functional.

\begin{exercise}
  Track the evolution of the letter \textit{R} from Phoenician
  \textit{resh} through Greek $\rho$ (rho) to Latin $R$.
  At each stage, identify the constraint pressure that drove the change.
  Express each transformation as a modification of the distinguishing
  feature set $F(g)$.
\end{exercise}

\begin{exercise}
  Korean hangul and the Cherokee syllabary were both designed to
  maximize learnability for a specific community. Compare their
  design strategies. What different constraints were the designers
  optimizing for? Which is more learnable for a native speaker of each
  language, and why?
\end{exercise}

\begin{exercise}[title={Design}]
  Design an optimal writing system for a language with exactly 20
  consonants and 5 vowels, using the constraint satisfaction framework.
  Your alphabet should minimize $\kappa_\text{total}$ with weights
  $w_1 = w_2 = w_3 = w_4 = 1/4$. What does your alphabet look like?
\end{exercise}
